\chapter{Coherence Clusters and Their Nerve}
\label{sec:clusters}
% ===================================================================

\section{Coherence Clusters}

\begin{definition}[Coherence Cluster]
A \emph{coherence cluster} is a maximal compact subpopulation:
a compact population $\mathcal{C} \subseteq \pop$ such that
for any agent $A \in \pop \setminus \mathcal{C}$, the population
$\mathcal{C} \cup \{A\}$ is not compact.
\end{definition}

\begin{remark}
Maximality means the cluster cannot absorb any additional agent
without losing its consensus property.
The boundary of a coherence cluster is the locus where one more
agent would introduce an irresolvable oscillation.
This is the precise topological meaning of the intuitive notion
of a \emph{coherence horizon}.
\end{remark}

\begin{remark}[Connection to Communication Radius]
The coherence clusters are the topological refinement of the
communication-radius clusters from Section~\ref{sec:population}.
The communication radius $B_\epsilon(A)$ gives a metric notion of
cluster; compactness gives a topological (and therefore
coordinate-free) notion.
The relationship between these two notions --- and in particular
whether they coincide under reasonable conditions on the GSLT ---
is an open question noted in Section~\ref{sec:gaps}.
\end{remark}

\section{The Nerve of the Covering}

\begin{definition}[Cluster Covering]
Let $\{\mathcal{C}_i\}_{i \in I}$ be the collection of all coherence
clusters of $\pop$.
This collection is a \emph{covering} of $\pop$: every agent belongs
to at least one coherence cluster.
\end{definition}

\begin{definition}[Nerve]
The \emph{nerve} $\nerve$ of the cluster covering is the simplicial
complex whose:
\begin{itemize}[itemsep=2pt]
  \item vertices are the coherence clusters $\mathcal{C}_i$;
  \item $k$-simplices are collections $\{\mathcal{C}_{i_0}, \ldots,
        \mathcal{C}_{i_k}\}$ of $k+1$ clusters with non-empty
        mutual intersection:
        $\mathcal{C}_{i_0} \cap \cdots \cap \mathcal{C}_{i_k} \neq \emptyset$.
\end{itemize}
\end{definition}

\begin{remark}[The Nerve as Effective Ontology]
The nerve $\nerve$ encodes the \emph{global topology} of the
population at the cluster scale.
Its vertices are the coherent communities; its edges are the
communities with common ground; its triangles are the communities
that all share a common ground; and so on.
This combinatorial object is the coarse-grained map of the
population's ontological landscape.
\end{remark}

\begin{remark}[Non-Overlapping Clusters]
Two clusters with empty intersection share no common ground.
No consensus is possible between them, even in principle.
They are \emph{ontologically disjoint}: there is no HML formula
that both communities can agree on.
In the context of the hierarchical population, these are the
communities that have drifted so far apart in their world models
that communication between them is not merely difficult but
meaningless.
\end{remark}

\begin{remark}[Overlapping Clusters]
Two clusters with non-empty intersection share a compact subspace
on which they agree.
This overlap is the fragment of ontology that both communities
can coherently discuss.
Cross-cluster communication is possible, but only through the
shared vocabulary of the overlap region.
The overlap is the \emph{common ground} in the precise topological
sense.
\end{remark}

\section{The Hierarchy}

The degradation argument of Section~\ref{sec:population} applies
recursively.
The coherence clusters are themselves agents (by assumption on the
population scale).
Among these cluster-agents, communication degrades at the
inter-cluster scale.
Applying the same argument, the cluster-agents fragment into
clusters-of-clusters.
The process iterates.

\begin{definition}[Ontological Hierarchy]
The \emph{ontological hierarchy} of a massive population $\pop$
is the sequence of nerves
\[
  \nerve_0 \;\leftarrow\; \nerve_1 \;\leftarrow\; \nerve_2
  \;\leftarrow\; \cdots
\]
where $\nerve_0$ is the nerve of the base-level clusters,
$\nerve_1$ is the nerve of the cluster-level clusters, and so on.
The arrows are coarse-graining maps (simplicial maps that collapse
finer structure).
\end{definition}

\begin{remark}
The ontological hierarchy is the computational analogue of the
hierarchy of scales in physics:
quarks $\to$ hadrons $\to$ nuclei $\to$ atoms $\to$ molecules
$\to$ cells $\to$ organisms $\to$ planets $\to$ galaxies
$\to$ superclusters.
Each level appears self-contained to an observer at that level.
The coarse-graining maps are the analogues of effective field theory
truncations: they discard the fine-grained structure below the
observer's resolution.
\end{remark}

% ===================================================================
